\section{Introduction}

A phishing email that reaches a security analyst has usually passed several
initial spam filters. Often, the next security layer is a small language model
(an AI text classifier). The model makes a fast, automated decision: is this
message malicious (harmful) or benign (safe)? Small open models with 1 to 4
billion parameters are popular for this job. They run on a single standard GPU,
can be hosted privately inside an organization, and require no special training
examples~\cite{rojas2024zeroshot,koide2024chatspamdetector}. However, an attacker
can change their message to bypass this filter.

This paper tests how much small text edits (called perturbations) help an
attacker evade detection. We limit the attacker to simple changes that a human
reader would easily overlook. These changes include look-alike characters (such
as replacing Latin letters with similar Cyrillic letters), tiny typos, and simple
synonyms. The total changes are kept under a strict 15\% edit budget per
message. We measure how often these edits cause the model's decision (verdict)
to flip from \emph{malicious} to \emph{benign}.

\paragraph{Why the denominator is crucial.}
An attack success rate (ASR) is meaningless without knowing its denominator
(the total number of tested samples). In security research, two different
denominators are often used. The first counts all malicious messages. The
second counts only the messages that the clean model correctly identified as
malicious before the attack. The second definition is what truly matters for
evasion: if a model never caught a message in the first place, modifying that
message cannot count as a successful attack flip~\cite{morris2020textattack}.
Therefore, we only attack messages that the clean model originally flagged as
malicious. We re-test every original message during the attack process to
confirm this label, and we report the baseline catch rate alongside every ASR.

\paragraph{Contributions.}
\begin{enumerate}
\item \textbf{A strict, reproducible evasion evaluation protocol}
  (\S\ref{sec:protocol}): We use a conditional denominator with re-verified
  original messages, an exact Levenshtein edit distance budget shared across all
  edits, identical prompts between baseline and attack phases, and public
  sample records so every result can be verified.
\item \textbf{An empirical evaluation of three small open models}
  (\S\ref{sec:results}): Testing 2{,}079 caught messages across 8{,}047 model
  queries resulted in 40 total decision flips. We test character-level and
  word-level edits under the same budget, provide 95\% confidence intervals for
  every row, and evaluate a saliency-guided attack that spends more queries to
  rank word importance (\S\ref{sec:saliency}).
\item \textbf{A paired comparison between edit types}
  (\S\ref{sec:pairs}): Word substitutions flip more decisions than
  character noise while requiring roughly half as many edits (26 versus 10
  sample-specific flips, $p=0.011$). In addition, ranking word swaps by model
  saliency flips more messages than random word swaps under the same budget
  (33 versus 7, $p=4\times10^{-5}$). Pairing tests on the exact same messages
  provides an honest comparison.
\item \textbf{Low transferability across models} (\S\ref{sec:transfer}):
  Replaying saved text edits across different models flips almost no samples off
  the diagonal. Transfer rates remain low overall. The edits appear specific to
  individual models rather than universal. Guided edits transfer slightly better
  than random edits, but remain ineffective as portable attacks.
\item \textbf{Analysis of model decision boundaries}: The model with the
  highest attack success rate is the one operating closest to its decision
  boundary, not the least accurate model overall. Conversely, a model that marks
  almost everything as malicious appears artificially robust simply because its
  classification rule rarely changes.
\item \textbf{A clear negative result}: Under a 15\% edit budget, random edits
  rarely fool these models. For practical deployments, the real concern is
  initial detection accuracy rather than vulnerability to simple random edits.
\end{enumerate}

\paragraph{Scope and limitations.} We do not evaluate white-box attacks with
full access to model internal weights, nor do we evaluate adaptive attackers who
know the defense beforehand. Our saliency search selects words using model
feedback, but it also uses a larger dictionary than the random word baseline.
Therefore, the gap between them is not due to word ordering alone
(\S\ref{sec:saliency}). An ablation study in \S\ref{sec:ablation} isolates these
factors and shows their individual impact on identical messages. Across the two
random attacks, we observe 40 flips in 2{,}079 messages across 8{,}047 queries.
Because these flip counts are small, the confidence intervals are wide, and we
focus on general trends rather than minor numerical differences.
